% !TEX root = ../../../main.tex
% Author-requested overview of number systems, containment, and cardinality.
\begin{pendingreview}
\section{Number systems and cardinalities}
\label{sec:analysis-i:number-systems}

\newthought{The roots just constructed are part of a larger picture.}
In \autoref{sec:analysis-i:roots-and-powers}, completeness supplied
$\sqrt2$ as a real number. Its equation $x^2-2=0$ also identifies it
as algebraic, in the sense of
\autoref{sec:analysis-i:algebraic-transcendental}. We now bring that
classification together with the chapter's counting results. The two
questions are where a number belongs and how large each set is.

The familiar number systems satisfy
\[
  \N\subsetneq\Q\subsetneq\R\subsetneq\C.
\]
Here $\N=\{1,2,3,\ldots\}$, $\Q$ consists of ratios of integers with
nonzero denominator, and $\R$ is the real line. A complex number is
$a+bi$ with $a,b\in\R$ and $i^2=-1$; we identify $a$ with $a+0i$.
The examples $1/2$, $\sqrt2$, and $i$ show, respectively, that each
inclusion is strict.

\subsection{Where algebraic and transcendental numbers belong}
Recall the two sets from
\autoref{sec:analysis-i:algebraic-transcendental}. The set
$\mathcal A$ consists of the real algebraic numbers, and
$\mathcal T=\R\setminus\mathcal A$ consists of the transcendental reals.
Inserting the algebraic numbers into the chain gives
\[
  \N\subsetneq\Q\subsetneq\mathcal A\subsetneq\R\subsetneq\C.
\]
Every rational is algebraic, while roots such as $\sqrt2$ show that
$\mathcal A$ also contains irrational numbers. The transcendental
reals include $\pi$ and $e$, whose transcendence was stated earlier
without proof. Thus the irrational reals split into
$\mathcal A\setminus\Q$ and $\mathcal T$.
Complex numbers can also be algebraic; $i$ satisfies $x^2+1=0$, but
is outside the \emph{real} algebraic set $\mathcal A$.
\autoref{fig:number-systems} shows these relationships with nested boxes.
Inside the real box, every number is either in the algebraic box or in
the transcendental region outside it.

\begin{figure}[htbp]
  \centering
  % Nested rectangles distinguish containment from the size of infinite sets.
\begin{tikzpicture}[analysis diagram]
  \draw (0,0) rectangle (10,7.2);
  \node[analysis label,anchor=north west] at (0.3,6.95)
    {$\C$: complex numbers};
  \node[analysis label,anchor=north east] at (9.65,6.95) {$i,\ 1+i$};

  \draw (0.35,0.35) rectangle (9.65,6.25);
  \node[analysis label,anchor=north west] at (0.65,6)
    {$\R$: real numbers};

  \draw (0.7,0.7) rectangle (6.3,5.3);
  \node[analysis label,anchor=north west] at (1,5)
    {$\mathcal A$: real algebraic numbers};
  \node[analysis label,anchor=north west] at (1,4.35) {$\sqrt2,\ \sqrt[3]{5}$};

  \draw (1.05,1.05) rectangle (5.95,3.7);
  \node[analysis label,anchor=north west] at (1.35,3.45)
    {$\Q$: rational numbers};
  \node[analysis label,anchor=east] at (5.6,2.5) {$\frac12,\ -\frac34$};

  \draw (1.4,1.4) rectangle (4.25,2.8);
  \node[analysis label,anchor=north west] at (1.7,2.55)
    {$\N$: naturals};
  \node[analysis label,anchor=north west] at (1.7,1.95)
    {$1,2,3,\ldots$};

  \node[analysis label,align=center] at (7.95,4.65)
    {Transcendental reals\\$\mathcal T=\R\setminus\mathcal A$};
  \node[analysis label] at (7.95,3.6) {$\pi,\ e$};
\end{tikzpicture}

  \caption[Containment of the number systems.]{Nested boxes show
    containment. Examples lie in the regions they belong to, outside
    the smaller boxes that exclude them. Areas are schematic and do
    not represent cardinalities.}
  \label{fig:number-systems}
\end{figure}

\subsection{Their infinite sizes}
Cardinality compares sets by bijections, rather than by containment.
The earlier algebraic/transcendental distinction also determines two
different infinite sizes. Combining the counting results in
\autoref{sec:analysis-i:sets-and-countable-sets} with
\autoref{thm:analysis-i:algebraic-transcendental-cardinalities} gives
\[
  |\N|=|\Q|=|\mathcal A|=\aleph_0,
  \qquad
  |\R|=|\mathcal T|=\mathfrak c=2^{\aleph_0}.
\]
Here $\aleph_0$, read ``aleph zero,'' is the smallest infinite
cardinality, and the continuum $\mathfrak c$ is strictly larger.
Thus the step from $\Q$ to all real algebraic numbers adds numbers
such as $\sqrt2$ without increasing cardinality. The transcendental
reals alone already have the cardinality of the whole real line.

Since $a+bi$ corresponds to the pair $(a,b)$, the square-cardinality
result in \autoref{ex:hw01:square-cardinality} and a bijection
$(0,1)\to\R$ give $|\C|=|\R^2|=\mathfrak c$.
A proper subset can therefore have the same cardinality as its
containing set.

\subsection{Aleph one and beyond}
The next infinite cardinalities are written
\[
  \aleph_0<\aleph_1<\aleph_2<\aleph_3<\cdots.
\]
The cardinal $\aleph_1$ is the smallest uncountable cardinality.
Each $\aleph_{n+1}$ is the next cardinality after $\aleph_n$, with
none strictly between them. The hierarchy continues beyond finite
subscripts; there is no largest infinite cardinality.

The \emph{Continuum Hypothesis} (CH) asserts that
$\mathfrak c=\aleph_1$. Uncountability alone does not establish this
equality. The bare $\aleph$ used in some exercises denotes
$\mathfrak c$.
\end{pendingreview}
